\section{Conclusiones}

\subsection{Sostenibilidad y escalabilidad del proyecto}

La implementación del ecosistema Odoo establece una hoja de ruta sólida que permite a Machu Picchu Exclusive Tours E.I.R.L. evolucionar desde una gestión manual y fragmentada hacia un modelo digital inteligente. Al centralizar la data operativa en la nube, la empresa cuenta ahora con una infraestructura digital robusta y escalable, capaz de manejar un volumen creciente de pasajeros en itinerarios complejos (de 2 a 20 días) sin incrementar la carga operativa de la gerencia. Esta transición no solo elimina la duplicidad de tareas, sino que profesionaliza la atención al cliente y facilita la futura expansión hacia nuevos mercados y socios internacionales.

\subsection{Mejora en la eficiencia y competitividad}

De acuerdo con el Plan de Implementación, la digitalización del ciclo comercial y la adopción de un CRM integral sientan las bases para mejorar de manera significativa la eficiencia operativa y la competitividad de la empresa en los mercados internacionales. Los beneficios esperados incluyen reducciones sustantivas en los tiempos de respuesta, disminución de errores administrativos asociados a la gestión manual y un aumento progresivo en la captación y fidelización de clientes a través de canales digitales directos.

Estas mejoras se proyectan como metas de mediano plazo y deberán ser verificadas durante la fase de monitoreo y evaluación del proyecto, utilizando los indicadores definidos (ventas digitales, tiempos de respuesta, recurrencia de clientes y reducción de errores), tal como se describe en el Plan.

\subsection{Fortalecimiento de la seguridad y reputación digital}

La implementación del repositorio estándar de Odoo para la gestión de documentos de identidad y pasaportes permite avanzar hacia un manejo más seguro de la información confidencial de los turistas. Al sustituir progresivamente el envío y almacenamiento volátil en dispositivos personales por un esquema centralizado y controlado, la empresa reduce los riesgos de pérdida o filtración de datos y fortalece las condiciones para consolidar una reputación digital sólida ante el mercado de lujo internacional.

\subsection{Trazabilidad y control operativo integral}

La integración de los módulos de CRM y Proyectos permite una visión 360° del turista, conectando el proceso de venta con la ejecución del servicio. El sistema garantiza la trazabilidad total de los itinerarios, permitiendo un control riguroso de los saldos pendientes y la carga de documentos operativos (vouchers y boletos). Este orden logístico es fundamental para asegurar el éxito de los tours y evitar errores en la compra de servicios no reembolsables, garantizando una operación profesional y confiable.

